\section{The impact of the new collider data}

\label{sec:impactnewdata}

We now study the dependence of the NNPDF3.1 PDF set upon 
the experimental information on which it is based. Firstly we disentangle
the effects of new data from the effects of methodological changes. Then 
we systematically quantify the impact on PDFs of each
new piece of experimental information added in
NNPDF3.1. Finally we discuss PDF determinations based on particular data subsets; 
PDFs determined only from collider data (i.e. excluding all fixed target data), 
 only from proton data (i.e. excluding all nuclear data), or excluding all LHC data.
As these PDF sets based on reduced dataset can also be useful for
specific phenomenological applications,  they are also made available
(see Section~\ref{sec:delivery} below). As in the previous Section, here
we will only present a selection of representative plots, the
interested reader is referred to a much larger set of plots available online 
as discussed in Section~\ref{sec:delivery}.


\subsection{Disentangling the effect of new data and methodology}
\label{sec:disentangling}

In Section~\ref{sec:results-mc} we have studied the impact of the main
methodological improvement introduced in NNPDF3.1, namely,
independently parametrizing the
charm PDF and determining it from the data. 
In order to completely disentangle the effect of data and methodology
we have performed a PDF determination using NNPDF3.1 methodology,
but the NNPDF3.0 dataset: specifically, we have removed from the
NNPDF3.1 dataset all the new data. There remain some small residual differences
between this restricted dataset and that of NNPDF3.0, specifically in 
some small differences in cuts and in the use of the combined HERA data
instead of the separate HERA-I and HERA-II sets. However, these differences
are expected to be minor~\cite{Rojo:2015nxa}.


%%%%%%%%%%%%%%%%%%%%%%%%%%%%%%%%%%%%%%%%%%%%%%%%%%%%%%%%%%%%%%%%%%%%%
\begin{figure}[t]
\begin{center}
  \includegraphics[scale=1]{images/plots/distances_nnpdf31_newdata.pdf}
  \caption{\small Same as Fig.~\ref{fig:distances_31_vs_30}, but now
    comparing the NNPDF3.1 NNLO
    global PDFs to PDFs determined using exactly the same methodology
    but with the NNPDF3.0 dataset.
    \label{fig:distances_nnpdf31_newdata}
  }
\end{center}
\end{figure}
%%%%%%%%%%%%%%%%%%%%%%%%%%%%%%%%%%%%%%%%%%%%%%%%%%%%%%%%%%%%%%%%%%%%%%

In Fig.~\ref{fig:distances_nnpdf31_newdata} we show the
distances between the NNPDF3.1 NNLO PDF set,  and that 
based on the NNPDF3.0 dataset using the same methodology.
%
We see that the impact of the new data is mostly localized at
large $x$, for the up, down and charm quarks and the gluon,
and at medium $x$ for strangeness.
%
As far as uncertainties are concerned, we observe improvements of
up to half a sigma across a wide range in $x$ and for all PDF flavors.
In Fig.~\ref{fig:31-nnlo-old-vs-new} we compare some representative PDFs for 
NNPDF3.1, the set based on NNPDF3.0 data with NNPDF3.1 methodology, and the 
original NNPDF3.0. We
see that the overall effect of the new data and the new methodology are 
comparable, but that they act in different regions and for different PDFs.
     For instance, for the light quarks and the gluon the impact of the
     new methodology dominates for all $x\lsim 10^{-2}$, where it
     produces an enhancement, and specifically the enhancement of the
     gluon for $x\lsim 0.03$ which was discussed in Sect.~\ref{sec:PDFcomparisons}.
      At large $x$ instead the dominant effect is
     from the new data, which lead to a reduction of the gluon
     and an enhancement of the quarks. Whereas of course charm is very
     significantly affected by the change in methodology --- it was
     not independently parametrized
     in NNPDF3.0 --- for $x\gsim0.1$, the new data also have a big
     impact. In fact, while strangeness is mostly affected
     by the new data in the medium and small $x$ regions, charm and gluon are  
      most affected by them at large $x$.

%%%%%%%%%%%%%%%%%%%%%%%%%%%%%%%%%%%%%%%%%%%%%%%%%%%%%%%%%%%%%%%%%%%%%
\begin{figure}[t]
\begin{center}
  \includegraphics[scale=0.38]{images/plots/xg-31-nnlo-old-vs-new.pdf}
  \includegraphics[scale=0.38]{images/plots/xu-31-nnlo-old-vs-new.pdf}
  \includegraphics[scale=0.38]{images/plots/xd-31-nnlo-old-vs-new.pdf}
  \includegraphics[scale=0.38]{images/plots/xdbar-31-nnlo-old-vs-new.pdf}
   \includegraphics[scale=0.38]{images/plots/xsp-31-nnlo-old-vs-new.pdf}
   \includegraphics[scale=0.38]{images/plots/xc-31-nnlo-old-vs-new.pdf}
   \caption{\small Same as Fig.~\ref{fig:31-nnlo-vs30}, but now also
     including PDFs determined using NNPDF3.1 methodology with the NNPDF3.0
     dataset. From left to right and from top to bottom the gluon, up,
     down, antidown,
     total strangeness and charm are shown.
    \label{fig:31-nnlo-old-vs-new}
  }
\end{center}
\end{figure}
%%%%%%%%%%%%%%%%%%%%%%%%%%%%%%%%%%%%%%%%%%%%%%%%%%%%%%%%%%%%%%%%%%%

\subsection{The transverse momentum of the $Z$ boson}
\label{sec:impactzpt}
%
The use of transverse momentum distributions has been advocated for a
long time (see e.g. Ref.~\cite{Forte:2013wc}) as a clean and powerful
constraint on PDFs, particularly the gluon.
As discussed in Section~\ref{sec:expdata}, it is now 
possible to include such data at NNLO thanks to the availability of the 
computation of this process up to NNLO QCD, along with precise data on $Z$ $p_T$ 
from ATLAS and CMS at 8 TeV. The impact of this dataset on PDFS  
has recently been studied in detail in Ref.~\cite{Boughezal:2017nla}.

 NNPDF3.1 is the first global PDF determination to include this data. 
In order to assess the impact of this dataset, we have repeated the NNLO
determination, excluding all $Z$ $p_T$ data. 
In Fig.~\ref{fig:distances_noZpT} we show the
distances between this PDF set and the default: it is clear
that the effect on all PDFs is moderate, with changes below one third of a sigma. 
The largest differences are seen in the gluon, as
expected, and the strange distributions. The reason for this state of affairs can
be best understood by directly comparing PDFs and their uncertainties,
see Fig.~\ref{fig:31-nnlo-ZpT}. It is clear that central values move
very little while uncertainties are slightly reduced, therefore
demonstrating the excellent consistency of the constraint from these measurements
with the existing dataset. The $Z$ $p_T$  dataset therefore reinforces the reliability of our gluon
determination. It also reduces somewhat the uncertainty on the total strangeness.
%
While in Ref.~\cite{Boughezal:2017nla} this dataset was found to have a rather
stronger impact than shown here, it should be noted that this was the
case when
determining PDFs from  the NNPDF3.0 dataset (less
 the jet data). In NNPDF3.1 more data are added, specifically top pair
 differential distributions:  a smaller impact
 of the $Z$ $p_T$ dataset when added to a wider prior  is not unexpected.

%%%%%%%%%%%%%%%%%%%%%%%%%%%%%%%%%%%%%%%%%%%%%%%%%%%%%%%%%%%%%%%%%%%%%
\begin{figure}[t]
\begin{center}
  \includegraphics[scale=1]{images/plots/distances_noZpT.pdf}
  \caption{\small Same as Fig.~\ref{fig:distances_31_vs_30},  but now
    comparing the default NNPDF3.1 to a version of it with the  
8 TeV $Z$ $p_T$ data from ATLAS and CMS not included.
\label{fig:distances_noZpT}}
\end{center}
\end{figure}
%%%%%%%%%%%%%%%%%%%%%%%%%%%%%%%%%%%%%%%%%%%%%%%%%%%%%%%%%%%%%%%%%%%%%%


%%%%%%%%%%%%%%%%%%%%%%%%%%%%%%%%%%%%%%%%%%%%%%%%%%%%%%%%%%%%%%%%%%%%%
\begin{figure}[t]
\begin{center}
  \includegraphics[scale=0.38]{images/plots/xg-31-nnlo-ZpT.pdf}
  \includegraphics[scale=0.38]{images/plots/xsp-31-nnlo-ZpT.pdf}
  \includegraphics[scale=0.38]{images/plots/xg-ERR-31-nnlo-ZpT.pdf}
  \includegraphics[scale=0.38]{images/plots/xsp-ERR-31-nnlo-ZpT.pdf}
  \caption{\small Same as Fig.~\ref{fig:31-nnlo-vs30} (top) and as
    Fig.~\ref{fig:ERR-31-nnlo-vs30} (bottom), but now
    comparing the default NNPDF3.1 to a version of it with the  
8 TeV $Z$ $p_T$ data from ATLAS and CMS not included. Results are
shown for the gluon (left) and total strangeness (right).
\label{fig:31-nnlo-ZpT}}
\end{center}
\vskip-0.5cm
\end{figure}
%%%%%%%%%%%%%%%%%%%%%%%%%%%%%%%%%%%%%%%%%%%%%%%%%%%%%%%%%%%%%%%%%%%%%%

In addition to the 8 TeV measurements from ATLAS and CMS there also exists 
a measurement of the normalized distribution at 7 TeV from ATLAS. The inclusion
of this dataset is problematic because the covariance matrix for a
normalized distribution depends on the cuts imposed on the dataset,
and only the covariance matrix for the full dataset is available. This
issue was studied in detail in
Ref.~\cite{Boughezal:2017nla}. Furthermore, this dataset is
superseded by the more precise 8~TeV measurement. Therefore, it
has not been included in the NNPDF3.1 dataset. However, we have
studied its potential impact by including it in a dedicated PDF determination, 
with its nominal published covariance matrix unmodified
despite the cuts.

In Table~\ref{tab:NNLOZpTchi2} we provide
 the $\chi^2/N_{\rm dat}$ values for all the LHC
    $Z$ $p_T$ measurements, for both the NNPDF3.1 NNLO baseline 
(including  the ATLAS and CMS $Z$ $p_T$ 8 TeV data), and also from 
the determination also including the ATLAS $Z$ $p_T$ 7 TeV data.
    %
    In the first column, the value of the $\chi^2/N_{\rm dat}$ for the 7 TeV data
    is in parenthesis to indicate that, unlike all other values, 
   it is a prediction
    and not the outcome of a fit. It is clear that
    the ATLAS $Z$ $p_T$ 7~TeV dataset is very poorly reproduced by the default NNPDF3.1
    set, and even after its inclusion in the dataset it cannot be
    accommodated. In fact, its inclusion is accompanied by a
    deterioration in the fit quality to ATLAS 8~TeV data, which are
    more accurate and supersede them. Furthermore, 
    there are also indications of tension between this dataset and the ATLAS
    $W$/$Z$ rapidity 
    distributions, whose total $\chi^2$ deteriorates by
    $12$ units (with 46 datapoints). 
The distances between these two PDF sets, displayed 
in Fig.~\ref{fig:distances_effect_Zpt7}, show that 
the gluon and quarks are shifted by almost one sigma by the
inclusion of the ATLAS $Z$ $p_T$ 7~TeV data. This is explicitly
shown in Fig.~\ref{fig:31-nnlo-AZPT7TEV} for the gluon and down
quark. It is apparent that uncertainties are however almost unchanged by the
inclusion of this dataset. 


%%%%%%%%%%%%%%%%%%%%%%%%%%%%%%%%%%%%%%%%%%%%%%%%%%%%%%%%%%%%%%%%%%%%%%
\begin{table}[b!]
  \centering
  \small
  \begin{tabular}{|c|c|c|}
    \hline
   &         NNPDF3.1 NNLO   &  + ATLAS $Z$ $p_T$ 7 TeV data \\
    \hline
    \hline
ATLAS $Z$ $p_T$ 7 TeV $(p_T^{ll},y_{ll})$    &     [6.78]           &   3.40    \\
ATLAS $Z$ $p_T$ 8 TeV $(p_T^{ll},M_{ll})$    &       0.93         &    0.98   \\
ATLAS $Z$ $p_T$ 8 TeV $(p_T^{ll},y_{ll})$     &    0.93            &   1.17    \\
CMS $Z$ $p_T$ 8 TeV $(p_T^{ll},M_{ll})$     &      1.32          &   1.33    \\
\hline
  \end{tabular}
  \caption{\small
    The values of $\chi^2/N_{\rm dat}$ for the LHC
    $Z$ $p_T$ data using the NNPDF3.1 NNLO PDF set, and for a new PDF
    determination which
    also includes the ATLAS $Z$ $p_T$ 7 TeV data.
    \label{tab:NNLOZpTchi2}
  }
\end{table}
%%%%%%%%%%%%%%%%%%%%%%%%%%%%%%%%%%%%%%%%%%%%%%%%%%%%%%%%%%%%%%%%%%%%%%


While we cannot say how a better treatment of the
covariance matrix would affect the results, we must conclude that
within our current level of understanding, 
inclusion of the ATLAS 7~TeV $Z$ $p_T$ dataset would have a significant
impact on PDFs, without an improvement in precision, and with signs of
tension between this dataset and both the remaining $Z$ $p_T$ datasets, and
other $W$ and $Z$ production data. Therefore its inclusion in the global dataset
does not appear to be justified.

%%%%%%%%%%%%%%%%%%%%%%%%%%%%%%%%%%%%%%%%%%%%%%%%%%%%%%%%%%%%%%%%%%%%%
\begin{figure}[t]
\begin{center}
  \includegraphics[scale=1]{images/plots/distances_effect_Zpt7.pdf}
  \caption{\small Same as Fig.~\ref{fig:distances_31_vs_30},  but now
    comparing the default NNPDF3.1 to a version of it with  
the 7~TeV $Z$ $p_T$ ATLAS data also included.
 \label{fig:distances_effect_Zpt7}}
\end{center}
\end{figure}
%%%%%%%%%%%%%%%%%%%%%%%%%%%%%%%%%%%%%%%%%%%%%%%%%%%%%%%%%%%%%%%%%%%%%%

%%%%%%%%%%%%%%%%%%%%%%%%%%%%%%%%%%%%%%%%%%%%%%%%%%%%%%%%%%%%%%%%%%%%%
\begin{figure}[t]
\begin{center}
  \includegraphics[scale=0.38]{images/plots/xg-31-nnlo-AZPT7TEV.pdf}
  \includegraphics[scale=0.38]{images/plots/xd-31-nnlo-AZPT7TEV.pdf}
  \caption{\small 
Same as Fig.~\ref{fig:31-nnlo-vs30}  but now
    comparing the default NNPDF3.1 to a version of it  with  
the 7~TeV $Z$ $p_T$ ATLAS data also included. Results are
shown for the gluon (left) and down quark (right).
\label{fig:31-nnlo-AZPT7TEV}}
\end{center}
\end{figure}
%%%%%%%%%%%%%%%%%%%%%%%%%%%%%%%%%%%%%%%%%%%%%%%%%%%%%%%%%%%%%%%%%%%%%%

\subsection{Differential distributions for top  pair production}
\label{sec:impacttop}

The impact of differential top pair production on PDFs
and the optimal selection of top datasets has been discussed
extensively in Ref.~\cite{Czakon:2016olj}.
Here we briefly study the impact of the top data on NNPDF3.1 by
comparing with PDFs determined removing the top data from the dataset.
In Fig.~\ref{fig:distances_notop} we show the distances between these
PDF sets. Large differences can be seen in the gluon central value and uncertainty
for $x\gsim 0.1$: 
these data constrain the gluon for values as large as $x\simeq 0.6$~\cite{Czakon:2016olj}, a region in which
constraints from other processes are not available. The effect on other PDFs is moderate, with the largest
impact seen on charm at small $x$.

%%%%%%%%%%%%%%%%%%%%%%%%%%%%%%%%%%%%%%%%%%%%%%%%%%%%%%%%%%%%%%%%%%%%%
\begin{figure}[t]
\begin{center}
  \includegraphics[scale=1]{images/plots/distances_notop.pdf}
  \caption{\small 
Same as Fig.~\ref{fig:distances_noZpT} but now excluding all top data (total cross-sections and differential
distributions). Note the different scale on the $y$ axis in the left plot.
    \label{fig:distances_notop}
  }
\end{center}
\end{figure}
%%%%%%%%%%%%%%%%%%%%%%%%%%%%%%%%%%%%%%%%%%%%%%%%%%%%%%%%%%%%%%%%%%%%%%
The differences between the two PDF sets are demonstrated in Fig.~\ref{fig:31-nnlo-top}, where the gluon and the
charm quark are shown. There is a substantial reduction in the uncertainty of the large $x$ gluon, 
with the central value without top data being considerably higher than 
 the narrow error band of the result when top is included. This
 suggests a significant increase in the precision of the gluon determination
 due to the top data. For the large $x$ gluon the differences between NNPDF3.1 and
NNPDF3.0 seen in Figs.~\ref{fig:31-nnlo-vs30}-\ref{fig:31-nnlo-old-vs-new} are 
therefore partly driven by the top data.
%
The impact on quark PDFs is marginal, as can be seen in the case of charm.

As already mentioned in Sect.~\ref{sec:topdata}, it has been shown in
Ref.~\cite{Czakon:2016vfr} that the sensitivity of the rapidity
distribution on the top mass is minimal. In fact, in 
Ref.~\cite{Czakon:2016olj} it was shown 
that if the top mass is varied by~1~GeV,
 NLO theoretical predictions for the normalized rapidity distributions
at the LHC~8~TeV 
vary by 0.6\% at most in the kinematic range covered by the data,
which is much less than the uncertainty  on the data, or the size of
the  NNLO corrections. This strongly suggests that our results are essentially
independent of the value of the top mass.

%%%%%%%%%%%%%%%%%%%%%%%%%%%%%%%%%%%%%%%%%%%%%%%%%%%%%%%%%%%%%%%%%%%%%
\begin{figure}[t]
\begin{center}
  \includegraphics[scale=0.38]{images/plots/xg-31-nnlo-top.pdf}
  \includegraphics[scale=0.38]{images/plots/xc-31-nnlo-top.pdf}
  \includegraphics[scale=0.38]{images/plots/xg-ERR-31-nnlo-top.pdf}
  \includegraphics[scale=0.38]{images/plots/xc-ERR-31-nnlo-top.pdf}
  \caption{\small 
Same as Fig.~\ref{fig:31-nnlo-ZpT} but now excluding all top data (total cross-sections and differential
distributions).
Results are shown for the gluon (left) and charm (right), the PDFs above and their uncertainties below.
    %
\label{fig:31-nnlo-top}}
\end{center}
\end{figure}
%%%%%%%%%%%%%%%%%%%%%%%%%%%%%%%%%%%%%%%%%%%%%%%%%%%%%%%%%%%%%%%%%%%%%%

%\clearpage

\subsection{Inclusive jet production}
\label{sec:jetdata}

While jet data have been used for PDF determination for a long time,
their full NNLO treatment is only becoming possible now, thanks to the recent
completion of the relevant computation~\cite{Currie:2016bfm,Currie:2017ctp}.
However,
as discussed in Section~\ref{sec:datajets}, NNLO
 corrections are not yet available for
all datasets included in NNPDF3.1. Consequently in the default
NNPDF3.1 PDF determination, jets have been included using NNLO PDF
      evolution and NLO matrix elements supplemented by an extra
      theory uncertainty determined through scale variation. Here we
      assess generally the effect of jet data, and in particular the
      possible impact of this approximation.

To this end, we first repeat the NNPDF3.1 determination
but excluding jet data. The distances between these PDFs and the
default are shown
in Fig.~\ref{fig:distances_nojets}. It is clear that 
jet data have a moderate and very localized impact, on the gluon
in the region $0.1\lsim x\lsim 0.6$, at most at the half-sigma level,
with essentially no impact on other PDFs. The changes in all other PDFs
are compatible with a statistical fluctuation.
A direct comparison of the gluon PDFs and their uncertainties in
Fig.~\ref{fig:jetdata-nnlo} confirms this. The uncertainty on the gluon is reduced by up to a factor of 
two by the jet data in this region, with the central value of the gluon within the narrower 
uncertainty band of the default set.

It is interesting to observe that in NNPDF3.0 the impact of the jet
data was rather more significant, with uncertainties being reduced by
a large factor for all $x\gsim 0.1$. In NNPDF3.1 the gluon at large $x$ is 
strongly constrained by the top data, as discussed in
Section~\ref{sec:impacttop}. Specifically, 
the addition of the jet data leaves the gluon unchanged in this
region, see Fig.~\ref{fig:jetdata-nnlo}, but addition of the top data produces a significant shift, as
seen in Fig.~\ref{fig:31-nnlo-top}. This 
suggests excellent compatibility between the jet and top data, with
the large $x$ gluon now mostly determined by the top data. This also
explains the insensitivity to the NNLO correction to jet production,
to be discussed shortly. 


%%%%%%%%%%%%%%%%%%%%%%%%%%%%%%%%%%%%%%%%%%%%%%%%%%%%%%%%%%%%%%%%%%%%%
\begin{figure}[t]
\begin{center}
  \includegraphics[scale=1]{images/plots/distances_nojets.pdf}
  \caption{\small  Same as Fig.~\ref{fig:distances_noZpT} but now excluding all jet data.
    \label{fig:distances_nojets}
  }
\end{center}
\end{figure}
%%%%%%%%%%%%%%%%%%%%%%%%%%%%%%%%%%%%%%%%%%%%%%%%%%%%%%%%%%%%%%%%%%%%%%

Despite there reduced impact, the jet data still
play a non-negligible role. One may therefore worry about the reliability of the
theoretical treatment, based on NLO matrix elements with theory
uncertainties. In order to assess this, we have repeated the PDF
determination but now
using full NNLO theory in the case of the 2011 
7~TeV LHC jet data where it is available.
      The other jet datasets, namely the CDF Run II $k_T$ jets, the
      ATLAS and CMS $\sqrt{s}=2.76$ TeV datasets, and the ATLAS 2010
      7~TeV dataset, are treated as in
      the baseline. Essentially no change in PDFs is found, as
      illustrated in Fig.~\ref{fig:jetdataex-nnlo} where the gluon
      and down PDFs are shown.
      Such a result is consistent with the percent level NNLO corrections
      found when using our choice of the jet $p_T$ as the central scale, shown in 
      Fig.~\ref{sec:datajets}. 

Also, as mentioned in
      Sect.~\ref{sec:datajets}, only the central rapidity bin of the
      ATLAS 2011 7~TeV data has been included, because we have found
      that, while a good description can be achieved if each of the
      rapidity bins is included in turn, or if the uncertainties are
      decorrelated between rapidity bins, it is impossible to achieve
      a good description of all rapidity bins with correlations
      included. One may therefore wonder whether the inclusion of
      other rapidity bins would lead to different results for the
      PDFs, despite the fact that they have less PDF
      sensitivity~\cite{Rojo:2014kta}.  In order to check this, we
      have compared to the data the prediction for all of the ATLAS 2011 7~TeV
      data using the default NNPDF3.1 set, and  determined the
      $\chi^2$ for each rapidity bin separately. For the five rapidity
      bins which have not been included, from central to forward, we
      find $\chi^2/N_{\rm dat}=1.27,\>0.95,\>1.06,\>0.97,\>0.73$, with
      respectively $N_{\rm dat}=29,\>26,\>23,\>19,\>12$, to be
      compared to the value  $\chi^2/N_{\rm dat}=1.06$ for $N_{\rm
        dat}=31$ of Tab~\ref{tab:NNLOjetschi2} for the central
      rapidity bin which is  included. We 
      conclude that all rapidity bins are well reproduced, and thus
      none of them can have a significant pull on PDFs that might
      change the result if they were included.

      

%%%%%%%%%%%%%%%%%%%%%%%%%%%%%%%%%%%%%%%%%%%%%%%%%%%%%%%%%%%%%%%%%%%%%
\begin{figure}[t]
\begin{center}
  \includegraphics[scale=0.38]{images/plots/xg-31-nnlo-jets.pdf}
  \includegraphics[scale=0.38]{images/plots/xg-ERR-31-nnlo-jets.pdf}
    \caption{\small Comparison between the default NNPDF3.1 NNLO PDFs
      an alternative determination in which all jet data have been
      removed: the gluon (left) and the percentage uncertainty on it
      (right) are shown
\label{fig:jetdata-nnlo}}
\end{center}
\end{figure}

%%%%%%%%%%%%%%%%%%%%%%%%%%%%%%%%%%%%%%%%%%%%%%%%%%%%%%%%%%%%%%%%%%%%%%



%%%%%%%%%%%%%%%%%%%%%%%%%%%%%%%%%%%%%%%%%%%%%%%%%%%%%%%%%%%%%%%%%%%%%%
\begin{figure}[t]
\begin{center}
  \includegraphics[scale=0.38]{images/plots/xg-31-nnlo-exactjets.pdf}
  \includegraphics[scale=0.38]{images/plots/xd-31-nnlo-exactjets.pdf}
    \caption{\small 
Same as Fig.~\ref{fig:31-nnlo-vs30}  but now
comparing the default NNPDF3.1 NNLO PDFs
      to an alternative determination in which 
ATLAS and CMS 7~TeV jet data have been included
      using exact NNLO theory. The gluon (left) and down
      (right) PDFs are shown.
\label{fig:jetdataex-nnlo}}
\end{center}
\end{figure}

%%%%%%%%%%%%%%%%%%%%%%%%%%%%%%%%%%%%%%%%%%%%%%%%%%%%%%%%%%%%%%%%%%%%%%

In order to understand better the impact of the NNLO corrections and
their effect on the PDFs, in 
Fig.~\ref{fig:jetdata-nnlo-datath} we compare the
best-fit prediction to 
the 7~TeV  2011 CMS 
    and ATLAS data for the two PDF sets compared in
    Figs.~\ref{fig:distances_nojets}-\ref{fig:jetdataex-nnlo}.
The corresponding values of $\chi^2/N_{\rm dat}$ are collected in
Table~\ref{tab:NNLOjetschi2}, both for these and all other jet data.
First of all, one should note that for all experiments for which an exact NNLO
 computation is not yet available, listed on the top part of
Table~\ref{tab:NNLOjetschi2}, the $\chi^2$ values obtained with these two PDF
sets are almost identical. Because the predictions for these experiments 
are computed using the same theory (NLO with extra scale uncertainty),
this shows that the change in PDFs is very small.
For the experiments for which NNLO theory is available, $\chi^2$
values also change very little. A comparison of the data to theory (see
Fig.~\ref{fig:jetdata-nnlo-datath}) shows that the NLO and NNLO
predictions are, as expected, quite close. Nevertheless, the NNLO
prediction is in slightly better agreement with the data, and this is
reflected by the
better value of the NNLO $\chi^2$ shown in
Table~\ref{tab:NNLOjetschi2}, despite the
the extra scale uncertainty (shown as an inner error bar on the data
in Fig.~\ref{fig:jetdata-nnlo-datath}) that is added to
the NLO prediction only.

As a final check, we have repeated the PDF determination but with a cut excluding
large $p_T$ jet data, for which the NNLO corrections are larger
(compare Fig.~\ref{fig:jetcf}). Specifically, we have only kept 7~TeV
jet data with $p_T\le 240$~GeV,  2.76~TeV
jet data with $p_T\le 95$~GeV,  1.96~TeV
jet data with $p_T\le 68$~GeV. We found that the ensuing PDFs are
statistically equivalent (distances of order one-two) to those from
the default set.  
We conclude that the impact of the approximate treatment of jet data
on the default NNPDF3.1 set is very small.







%%%%%%%%%%%%%%%%%%%%%%%%%%%%%%%%%%%%%%%%%%%%%%%%%%%%%%%%%%%%%%%%%%%%%%
\begin{table}[H]
  \centering
  \begin{tabular}{|c|c|c|}
    \hline
   &         NNPDF3.1   &  exact NNLO  \\
    \hline
    \hline
CDF Run II $k_t$ jets     &          0.84      &      0.85 \\
ATLAS jets 2.76 TeV   & 1.05       &     1.03 \\
CMS jets 2.76 TeV  & 1.04        &    1.02 \\
ATLAS jets 2010 7 TeV  & 0.96     &       0.95 \\
\hline
ATLAS jets 2011 7 TeV &  1.06      &      0.91 \\
CMS jets 7 TeV 2011 7 TeV &   0.84   &         0.79 \\
\hline
  \end{tabular}
  \caption{\small
    The values of $\chi^2/N_{\rm dat}$ for all the jet datasets obtained
    using either of the two PDF sets compared in
    Fig.~\ref{fig:jetdataex-nnlo}, and in each case the theory used in
    the corresponding PDF determination. For the datasets in the top part of the
    table the exact NNLO computations are not yet available and NLO theory
    with scale uncertainty is used throughout, while for those in the
    bottom part of the table NNLO theory is used for the right column.
    \label{tab:NNLOjetschi2}
  }
\end{table}
%%%%%%%%%%%%%%%%%%%%%%%%%%%%%%%%%%%%%%%%%%%%%%%%%%%%%%%%%%%%%%%%%%%%%%

\clearpage


%%%%%%%%%%%%%%%%%%%%%%%%%%%%%%%%%%%%%%%%%%%%%%%%%%%%%%%%%%%%%%%%%%%%%
\begin{figure}[t]
\begin{center}
  \includegraphics[width=0.49\textwidth]{images/plots/datath-cmsjets.pdf}
  \includegraphics[width=0.49\textwidth]{images/plots/datath-atlasjets.pdf}
  \caption{\small Comparison between CMS (left)
    and ATLAS (right) one-jet inclusive data 
    at 7~TeV from 2011, and best-fit results obtained using
NLO theory supplemented by scale uncertainties or exact NNLO
theory. The uncertainties shown on the best-fit prediction is the PDF
uncertainty, while that on the data is the diagonal (outer error bar)
and the scale uncertainty on the NLO prediction (inner error bar).
\label{fig:jetdata-nnlo-datath}}
\end{center}
\end{figure}
%%%%%%%%%%%%%%%%%%%%%%%%%%%%%%%%%%%%%%%%%%%%%%%%%%%%%%%%%%%%%%%%%%%%%%


\subsection{Electroweak boson production in the forward region}
\label{sec:lhcbimpact}

%%%%%%%%%%%%%%%%%%%%%%%%%%%%%%%%%%%%%%%%%%%%%%%%%%%%%%%%%%%%%%%%%%%%%
\begin{figure}[t]
\begin{center}
  \includegraphics[scale=1]{images/plots/distances_noLHCb.pdf}
  \caption{\small Same as Fig.~\ref{fig:distances_noZpT} but now
    excluding all LHCb data. Note the different scale on the $y$ axis 
in the left plot.
    \label{fig:distances_noLHCb}
  }
\end{center}
\end{figure}
%%%%%%%%%%%%%%%%%%%%%%%%%%%%%%%%%%%%%%%%%%%%%%%%%%%%%%%%%%%%%%%%%%%%%%

Electroweak production data from the LHCb experiment open up a new
kinematic region and therefore provide new constraints on flavor separation
at large and small $x$. While LHCb data were included in NNPDF3.0, 
the release of the legacy Run I LHCb measurements at 7 TeV and 8 TeV, which include all
correlations between $W$ and $Z$ data, greatly increases their utility in
PDF determination.
In Fig.~\ref{fig:distances_noLHCb} distances are shown
between the NNPDF3.1 NNLO default and PDFs determined excluding all
LHCb data.
%
The impact is significant for all quark PDFs, especially in the
valence region:
hence this data has a substantial impact on flavor
separation, most notably at large $x$.



%%%%%%%%%%%%%%%%%%%%%%%%%%%%%%%%%%%%%%%%%%%%%%%%%%%%%%%%%%%%%%%%%%%%%
\begin{figure}[t]
\begin{center}
  \includegraphics[scale=0.34]{images/plots/xu-31-nnlo-LHCb.pdf}
  \includegraphics[scale=0.34]{images/plots/xu-ERR-31-nnlo-LHCb.pdf}
  \includegraphics[scale=0.34]{images/plots/xd-31-nnlo-LHCb.pdf}
  \includegraphics[scale=0.34]{images/plots/xd-ERR-31-nnlo-LHCb.pdf}
  \includegraphics[scale=0.34]{images/plots/xc-31-nnlo-LHCb.pdf}
  \includegraphics[scale=0.34]{images/plots/xc-ERR-31-nnlo-LHCb.pdf}
  \caption{\small 
Same as Fig.~\ref{fig:31-nnlo-ZpT} but now excluding all LHCb data.
Results are presented, from top to bottow, for  the up, down and charm
PDFs. Both PDFs (left) and uncertainties (right) are shown.
\label{fig:lhcbdata-nnlo}}
\end{center}
\end{figure}
%%%%%%%%%%%%%%%%%%%%%%%%%%%%%%%%%%%%%%%%%%%%%%%%%%%%%%%%%%%%%%%%%%%%%%
This is explicitly demonstrated 
for the up, down and charm PDFs in
Fig.~\ref{fig:lhcbdata-nnlo}, where the percentage PDF 
uncertainties are 
also shown. The LHCb data play a
significant role in the data-driven large-$x$ enhancement of light quark
PDFs discussed in Section~\ref{sec:disentangling} (see
Fig.~\ref{fig:31-nnlo-old-vs-new}), and is largely responsible for
the sizable impact of new data on charm for $x\gsim0.1$, where they
significantly reduce the uncertainty.
 Effects are more marked at medium
and large $x$, peaking at around $x\simeq 0.3$: in this region the PDF
uncertainty is also substantially reduced; the reduction in
uncertainty is especially marked for the down PDF.

In order to see the impact of the LHCb data directly, in
Fig.~\ref{fig:lhcbdata-datatheory} we compare 
the 8~TeV LHCb muon $W^+$ and $W^-$  data
to predictions obtained using NNPDF3.0 and NNPDF3.1. The improvement is
clear, particularly for large rapidities. There is also a noticeable
reduction in PDF 
uncertainty on the prediction.

%%%%%%%%%%%%%%%%%%%%%%%%%%%%%%%%%%%%%%%%%%%%%%%%%%%%%%%%%%%%%%%%%%%%%
\begin{figure}[t]
\begin{center}
  \includegraphics[scale=0.45]{images/plots/LHCb_LHCBWZMU8TEV_dataset_report_datanorm_plot_fancy_0.pdf}
  \includegraphics[scale=0.45]{images/plots/LHCb_LHCBWZMU8TEV_dataset_report_datanorm_plot_fancy_1.pdf}
  \caption{\small Comparison between 8~TeV LHCb muon $W^+$ (left) and
    $W^-$ (right) production data to NNLO predictions obtained using
    NNPDF3.1 and NNPDF3.0. The uncertainties shown are the diagonal
experimental uncertainty for the data, and the PDF uncertainty for
the best-fit prediction. 
\label{fig:lhcbdata-datatheory}}
\end{center}
\end{figure}
%%%%%%%%%%%%%%%%%%%%%%%%%%%%%%%%%%%%%%%%%%%%%%%%%%%%%%%%%%%%%%%%%%%%%%

%%%%%%%%%%%%%%%%%%%%%%%%%%%%%%%%%%%%%%%%%%%%%%%%%%%%%%%%%%%%%%%%%%%%%
\begin{figure}[h]
\begin{center}
  \includegraphics[scale=1]{images/plots/distances_d0wasy.pdf}
  \caption{\small
Same as Fig.~\ref{fig:distances_noZpT} but now excluding 
 D0 $W$ asymmetry data.
    \label{fig:distances_nod0wasy}
  }
\end{center}
\end{figure}
%%%%%%%%%%%%%%%%%%%%%%%%%%%%%%%%%%%%%%%%%%%%%%%%%%%%%%%%%%%%%%%%%%%%%%



\subsection{$W$ asymmetries from the Tevatron}
\label{sec:legtevdata}

$W$ production data from the Tevatron have for many years been the
leading 
source of 
information on quark flavor decomposition. The final legacy 
D0 $W$ asymmetry measurements in the electron and muon
channels are included in NNPDF3.1, superseding all previous data.
In Fig.~\ref{fig:distances_nod0wasy} we perform a distance comparison between the default NNPDF3.1 and PDFs determined
excluding this dataset. Distances are generally small, an observation confirmed by direct PDF comparison in  Fig.~\ref{fig:d0wasydata-nnlo}.
However, we have seen in Tab.~\ref{tab:chi2tab_31-nlo-nnlo-30} that
the fit quality for this dataset is rather better with NNPDF3.1 than
with the  previous NNPDF3.0. The moderate impact of this dataset 
is due to its excellent consistency with the abundant LHC data, which are now driving flavor separation. 
This
data thus provides further evidence for the reliability
 of the flavor separation in NNPDF3.1.

%%%%%%%%%%%%%%%%%%%%%%%%%%%%%%%%%%%%%%%%%%%%%%%%%%%%%%%%%%%%%%%%%%%%%
\begin{figure}[h]
\begin{center}
  \includegraphics[scale=0.38]{images/plots/xubar-d0wasydata-nnlo.pdf}
  \includegraphics[scale=0.38]{images/plots/xdbar-d0wasydata-nnlo.pdf}
  \caption{\small 
Same as Fig.~\ref{fig:31-nnlo-ZpT} but now excluding D0 $W$
asymmetries.
The antiup (left) and antidown (right) PDFs are shown.
    \label{fig:d0wasydata-nnlo}
  }
\end{center}
\end{figure}
%%%%%%%%%%%%%%%%%%%%%%%%%%%%%%%%%%%%%%%%%%%%%%%%%%%%%%%%%%%%%%%%%%%%%

\subsection{The ATLAS $W,Z$ production data and strangeness}
\label{sec:atlaswz}

ATLAS $W$ and $Z$ production data were already included in NNPDF3.0,
but recent measurements based on the 
2011 dataset~\cite{Aaboud:2016btc} have much smaller statistical
uncertainties. This dataset, like the previous ATLAS measurement,
 has been claimed to have a large impact on
strangeness. This is borne out by the plot,
Fig.~\ref{fig:distances_atlaswz11rap}, of the distance between the
default NNPDF3.1 and a version from which this dataset has been
excluded. Indeed the largest effect --- almost at the one sigma level on central
values --- is seen on the strange and charm PDFs, with a rather smaller
impact on all other PDFs.

%%%%%%%%%%%%%%%%%%%%%%%%%%%%%%%%%%%%%%%%%%%%%%%%%%%%%%%%%%%%%%%%%%%%%
\begin{figure}[t]
\begin{center}
  \includegraphics[scale=1]{images/plots/distances_atlaswz11rap.pdf}
  \caption{\small
Same as Fig.~\ref{fig:distances_noZpT} but now excluding 
2011 ATLAS $W,Z$ rapidity
    distributions.
    \label{fig:distances_atlaswz11rap}
  }
\end{center}
\end{figure}
%%%%%%%%%%%%%%%%%%%%%%%%%%%%%%%%%%%%%%%%%%%%%%%%%%%%%%%%%%%%%%%%%%%%%%

A direct comparison of the strange and charm PDFs in
Fig.~\ref{fig:pdfs-noAWZrap11} shows that strangeness is significantly
enhanced in the medium/small $x$ region 
by the inclusion of the ATLAS data, while charm is
suppressed. As discussed in Section~\ref{sec:results-mc} and shown in
Fig.~\ref{fig:31-nnlo-fitted-vs-pch}, this suppression of charm cannot
be accommodated when charm is perturbatively generated, and therefore
parametrizing charm is important in order to be able to reconcile the ATLAS
data with the global dataset which in this $x$ range is severely constrained by
HERA data. The strange and charm content of the proton will be
discussed in detail in Sections~\ref{sec:strangeness} and \ref{sec:phenocharm} below.

%%%%%%%%%%%%%%%%%%%%%%%%%%%%%%%%%%%%%%%%%%%%%%%%%%%%%%%%%%%%%%%%%%%%%
\begin{figure}[t]
\begin{center}
  \includegraphics[scale=0.38]{images/plots/xsp-31-nnlo-noAWZrap11.pdf}
  \includegraphics[scale=0.38]{images/plots/xc-31-nnlo-noAWZrap11.pdf}
  \caption{\small 
Same as Fig.~\ref{fig:31-nnlo-ZpT} but now 
excluding 
2011 ATLAS $W,Z$ rapidity
    distributions.
The total strange (left) and charm (right) PDFs are shown.
    \label{fig:pdfs-noAWZrap11}
  }
\end{center}
\end{figure}
%%%%%%%%%%%%%%%%%%%%%%%%%%%%%%%%%%%%%%%%%%%%%%%%%%%%%%%%%%%%%%%%%%%%%


%%%%%%%%%%%%%%%%%%%%%%%%%%%%%%%%%%%%%%%%%%%%%%%%%%%%%%%%%%%%%%%%%%%%%
\begin{figure}[t]
\begin{center}
  \includegraphics[scale=0.42]{images/plots/ATLAS_ATLASWZRAP11_dataset_report_datanorm_plot_fancy_1.pdf}
  \includegraphics[scale=0.42]{images/plots/ATLAS_ATLASWZRAP11_dataset_report_datanorm_plot_fancy_2.pdf}
  \caption{\small Comparison between the 2011 ATLAS 7~TeV  $W^-$
    (left) and $Z$ (right)
    data to NNLO predictions obtained using NNPDF3.1 and NNPDF3.0; $W$
    production data 
     are plotted versus the pseudorapidity of the forward lepton
     $\eta_l$, while $Z$ production data are plotted vs the dilepton
     rapidity $y_{ll}$. 
 \label{fig:atlaswz11-datatheory}}
\end{center}
\end{figure}
%%%%%%%%%%%%%%%%%%%%%%%%%%%%%%%%%%%%%%%%%%%%%%%%%%%%%%%%%%%%%%%%%%%%%%

NNPDF3.1 achieves a good description of the data, as illustrated in 
Fig.~\ref{fig:atlaswz11-datatheory} where the NNLO prediction obtained
using NNPDF3.0 is also shown. It is clear that the agreement is
greatly improved, as demonstrated in the $\chi^2$ values shown
in Tab.~\ref{tab:chi2tab_31-nlo-nnlo-30}. It is also interesting to 
note the significant reduction in PDF uncertainties. As mentioned in
Sect.~\ref{data:inclusiveWZ}, these data have been included only
partially in the NNPDF3.1 determination: specifically $Z$ production
data off peak  or at forward rapidity have not been included. We have
checked however that the description of these data in NNPDF3.1 is
equally good, and similarly improved in comparison to NNPDF3.0,
with $\chi^2/N{\rm dat}$ of order unity. A comparison of these data
for two of the four bins which have not been included to
the NNPDF3.1 and NNPDF3.0 is shown in Fig.~\ref{fig:atlaswz11-datatheorypred}.

%%%%%%%%%%%%%%%%%%%%%%%%%%%%%%%%%%%%%%%%%%%%%%%%%%%%%%%%%%%%%%%%%%%%%
\begin{figure}[t]
\begin{center}
  \includegraphics[scale=0.42]{images/plots/ATLASWZRAP11CC_datanorm_highmasscentral.pdf}
  \includegraphics[scale=0.42]{images/plots/ATLASWZRAP11CF_datanorm_zpeakforward.pdf}
  \caption{\small Same as Fig.~\ref{fig:atlaswz11-datatheory} but now
    for two of the four data bins which have not been included in the
    NNPDF3.1 determination: high-mass $Z$ production at central rapidity (left)
    and on-shell $Z$ production at forward rapidity (right).
 \label{fig:atlaswz11-datatheorypred}}
\end{center}
\end{figure}
%%%%%%%%%%%%%%%%%%%%%%%%%%%%%%%%%%%%%%%%%%%%%%%%%%%%%%%%%%%%%%%%%%%%%%


\subsection{The CMS 8~TeV double-differential Drell-Yan distributions}
\label{sec:cms8tevimpact}

Like ATLAS, CMS has also published updated electroweak boson
production data. The NNPDF3.0 PDF determination already included 
double-differential (in rapidity and invariant mass) Drell-Yan data
at 7 TeV from the CMS 2011 dataset~\cite{Chatrchyan:2013tia}.
%
An updated version of the same measurement at 8~TeV based on 2012
data was presented in Ref.~\cite{CMS:2014jea}, including both the absolute 
cross-sections and the ratio of 8~TeV and 7~TeV measurements.

This data has very small uncorrelated systematic
uncertainties. Unfortunately, only the full covariance
matrix, with no breakdown of individual correlated systematics, has been made 
available. The combination of these two facts makes it impossible to include 
this experiment in the NNPDF3.1 dataset, as we now explain.
In Fig.~\ref{fig:distances_cmsdy2d12} we show
the distances between the NNPDF3.1 NNLO PDF set and a modified version of it
where this dataset has been included.
While the impact on uncertainties is moderate, clearly this dataset has a significant 
impact at the level of central-values on all PDFs for
almost all $x$ values, with a particularly important impact on 
the medium/small $x$ gluon. 
This is somewhat surprising, given that Drell-Yan
production only provides an indirect handle on the gluon PDF.

%%%%%%%%%%%%%%%%%%%%%%%%%%%%%%%%%%%%%%%%%%%%%%%%%%%%%%%%%%%%%%%%%%%%%
\begin{figure}[t]
\begin{center}
  \includegraphics[scale=1]{images/plots/distances_cmsdy2d12.pdf}
  \caption{\small
Same as Fig.~\ref{fig:distances_31_vs_30}  but now
    comparing the default NNPDF3.1 to a version of it  with  
the 8~TeV CMS double-differential Drell-Yan  data also included.
    \label{fig:distances_cmsdy2d12}
  }
\end{center}
\end{figure}
%%%%%%%%%%%%%%%%%%%%%%%%%%%%%%%%%%%%%%%%%%%%%%%%%%%%%%%%%%%%%%%%%%%%%%

A direct comparison of PDFs and their uncertainties 
in Fig.~\ref{fig:pdfs-cmsdy2d12} shows that these data induce an
upwards shift by up to one sigma of the gluon for $x\lsim 0.1$, and
a downward shift of the light quark PDFs for $x\gsim 0.1$, 
by a comparable amount. This, however, is not accompanied by a
reduction of PDF uncertainties, which increase a little, as
also shown in Fig.~\ref{fig:pdfs-cmsdy2d12}. 

Furthermore, while the fit quality of the 8~TeV CMS
double-differential Drell-Yan data remains poor after their inclusion in the fit,
with a value of $\chi^2/N_{\rm dat}=2.88$, there is a
certain deterioration in fit quality of all other experiments. Indeed,
the total $\chi^2$ to all the other data
deteriorates by $\Delta \chi^2=11.5$.
A more detailed inspection shows
that the most marked deterioration is seen in the HERA combined
inclusive DIS data, with $\Delta
\chi^2=19.7$. 
This means that there is tension between the CMS data and
the rest of the global dataset, and more specifically tension with the HERA data,
which are most sensitive to the small $x$ gluon.

We must conclude that this experiment appears to be inconsistent with
the global dataset, and particularly with the data 
with which we have the least reasons to doubt, namely the combined HERA data
and their determination of the gluon. In the absence of more detailed
information on the covariance matrix it is not possible to further
investigate the matter, and the 8~TeV CMS
double-differential Drell-Yan data have consequently not been included in the global dataset.

%%%%%%%%%%%%%%%%%%%%%%%%%%%%%%%%%%%%%%%%%%%%%%%%%%%%%%%%%%%%%%%%%%%%%
\begin{figure}[t]
\begin{center}
  \includegraphics[scale=0.38]{images/plots/xg-31-nnlo-cmsdy2d12.pdf}
  \includegraphics[scale=0.38]{images/plots/xu-31-nnlo-cmsdy2d12.pdf}
  \includegraphics[scale=0.38]{images/plots/xg-ERR-31-nnlo-cmsdy2d12.pdf}
  \includegraphics[scale=0.38]{images/plots/xu-ERR-31-nnlo-cmsdy2d12.pdf}
  \caption{\small 
Same as Fig.~\ref{fig:31-nnlo-vs30}  (top) but now
    comparing the default NNPDF3.1 to a version of it  with  
the 8~TeV CMS double-differential Drell-Yan  data also included. The
corresponding 
percentage uncertainties are also shown (bottom). Results are
shown for the gluon (left) and up quark (right).
    \label{fig:pdfs-cmsdy2d12}
  }
\end{center}
\end{figure}
%%%%%%%%%%%%%%%%%%%%%%%%%%%%%%%%%%%%%%%%%%%%%%%%%%%%%%%%%%%%%%%%%%%%%

\subsection{The EMC $F_2^c$ data and intrinsic charm.}
\label{sec:emc}

The advantages of introducing an independently parametrized charm PDF
were advocated in Ref.~\cite{Ball:2016neh}, where a first global PDF
determination including charm was presented, based on the NNPDF3.0 methodology and
dataset. The default NNPDF3.0 dataset was supplemented by charm
deep-inelastic structure function data from the EMC
collaboration~\cite{Aubert:1982tt}. This dataset is quite old, but it
remains the only measurement of the charm structure function in the
large $x$ 
region. With the wider NNPDF3.1 dataset the EMC dataset is no longer
quite so indispensable, specifically in view of phenomenology at the
LHC:
it has thus been omitted from the default NNPDF3.1 
determination as doubts have been raised about its reliability. 

However, a number of checks performed in
Refs.~\cite{Ball:2016neh,Rottoli:2016lsg},
such as variations of kinematical cuts and
systematic uncertainties, do not suggest any serious compatibility issues, and 
rather confirmed this dataset as being as reliable as the other older 
datasets with fixed nuclear targets routinely included in global PDF determinations.
Therefore it is interesting to revisit the issue of the impact of
this dataset within the context of 
NNPDF3.1. To this purpose we have produced a modified version of the global
NNPDF3.1  NNLO analysis in which the EMC dataset~\cite{Aubert:1982tt} is added to
the default NNPDF3.1 dataset. In Fig.~\ref{fig:distances_emc} 
the distances between this PDF set and the default are shown. 
The EMC dataset has a non-negligible impact on charm, at the one
sigma level, and also to a
lesser extent, at the half-sigma level,
on all light quarks, with only the gluon left essentially unaffected. 
The bulk of the effect is localized in the region   $0.01\lsim x\lsim 0.3$.
The PDFs are directly compared in Fig.~\ref{fig:31-nnlo-emc}.
The EMC data lead to an increase in the charm distribution
towards the upper edge of its error band in the default PDF set for
$0.02\lsim x\lsim 0.2$, while reducing the uncertainty
on it by a sizable factor. The light quark PDFs are correspondingly
slightly suppressed, and their uncertainties also reduced a little.


%%%%%%%%%%%%%%%%%%%%%%%%%%%%%%%%%%%%%%%%%%%%%%%%%%%%%%%%%%%%%%%%%%%%%
\begin{figure}[t]
\begin{center}
  \includegraphics[scale=1]{images/plots/distances_EMC.pdf}
  \caption{\small Same as Fig.~\ref{fig:distances_effect_Zpt7},  but now
    comparing the default NNPDF3.1 to a version of it with  
the EMC $F_2^c$ dataset also included.
    \label{fig:distances_emc}
  }
\end{center}
\end{figure}
%%%%%%%%%%%%%%%%%%%%%%%%%%%%%%%%%%%%%%%%%%%%%%%%%%%%%%%%%%%%%%%%%%%%%%



The values of $\chi^2/N_{\rm dat}$ for the
    deep-inelastic scattering experiments and for the total
    dataset before and after inclusion of the EMC data in the
    NNPDF3.1 dataset are shown in Table~\ref{tab:EMCchi2}. The fit quality of the 
    EMC charm dataset is greatly improved by its
    inclusion, without any significant change of the fit quality
    for any other DIS data: the hadron collider
    data are even less sensitive. The 
    inclusion of EMC charm data therefore appears to give a more accurate charm determination, 
    with no cost elsewhere, and so usage
    of this PDF set is recommended when precise charm PDFs at large $x$ are
    required. The phenomenological implications of the charm PDF will be
    discussed in Section~\ref{sec:phenocharm} below.

%%%%%%%%%%%%%%%%%%%%%%%%%%%%%%%%%%%%%%%%%%%%%%%%%%%%%%%%%%%%%%%%%%%%%%
\begin{table}[H]
  \centering
  \small
  \begin{tabular}{|c|c|c|}
    \hline
   &         NNPDF3.1 NNLO   &  + EMC charm data \\
    \hline
    \hline
NMC    &    1.30     &  1.29    \\
SLAC    &    0.75      &  0.76   \\
BCDMS    &    1.21      &  1.24  \\
CHORUS    &    1.11      & 1.10   \\
NuTeV dimuon    &    0.82      &   0.88    \\
\hline
HERA I+II inclusive    &    1.16      &    1.16       \\
HERA $\sigma_c^{\rm NC}$    &    1.45      &  1.42    \\
HERA $F_2^b$    &    1.11      &   1.11   \\
EMC $F_2^c$     &   [4.8]   &   0.93 \\
\hline
\hline
{\bf Total dataset}   &  {\bf  1.148}      &  {\bf  1.145}   \\
\hline
  \end{tabular}
  \caption{\small The values of $\chi^2/N_{\rm dat}$ for the
    deep-inelastic scattering experiments, as well as for the total
    dataset, for the
    NNPDF3.1 NNLO PDF set and for a new PDF determination which
    also includes the EMC charm structure function data.
    \label{tab:EMCchi2}
  }
\end{table}
%%%%%%%%%%%%%%%%%%%%%%%%%%%%%%%%%%%%%%%%%%%%%%%%%%%%%%%%%%%%%%%%%%%%%%

\clearpage


%%%%%%%%%%%%%%%%%%%%%%%%%%%%%%%%%%%%%%%%%%%%%%%%%%%%%%%%%%%%%%%%%%%%%
\begin{figure}[t]
\begin{center}
  \includegraphics[scale=0.38]{images/plots/xc-31-nnlo-emc.pdf}
  \includegraphics[scale=0.38]{images/plots/xc-ERR-31-nnlo-emc.pdf}
  \includegraphics[scale=0.38]{images/plots/xu-31-nnlo-emc.pdf}
  \includegraphics[scale=0.38]{images/plots/xd-31-nnlo-emc.pdf}
  \caption{\small 
Same as Fig.~\ref{fig:31-nnlo-AZPT7TEV} but now
    comparing the default NNPDF3.1 to a version of it with  
the EMC $F_2^c$ dataset also included. 
Results are
shown for the charm (top left),
up (bottom left) and
down (bottom right) PDFs. The  relative PDF
uncertainty on charm is also shown (top right).
    %
\label{fig:31-nnlo-emc}}
\end{center}
\end{figure}
%%%%%%%%%%%%%%%%%%%%%%%%%%%%%%%%%%%%%%%%%%%%%%%%%%%%%%%%%%%%%%%%%%%%%%


\subsection{The impact of LHC data}
\label{sec:nolhc}


We have seen that the LHC data have a significant impact on various
PDFs. Both in order to precisely gauge this impact, and in view of
possible applications in which usage of PDFs without LHC data is required,
we have produced a PDF set in which all LHC data are excluded from
the NNPDF3.1 dataset. The distance between the ensuing PDF set and the
default are shown in Fig.~\ref{fig:distances_lhc}. 

%%%%%%%%%%%%%%%%%%%%%%%%%%%%%%%%%%%%%%%%%%%%%%%%%%%%%%%%%%%%%%%%%%%%%
\begin{figure}[t]
\begin{center}
  \includegraphics[scale=1]{images/plots/distances_LHC.pdf}
  \caption{\small Same as Fig.~\ref{fig:distances_noZpT} but now
    excluding all LHC data. 
    \label{fig:distances_lhc}
  }
\end{center}
\end{figure}
%%%%%%%%%%%%%%%%%%%%%%%%%%%%%%%%%%%%%%%%%%%%%%%%%%%%%%%%%%%%%%%%%%%%%%

The cumulative effect of the data which were discussed in
Sects.~\ref{sec:impactzpt}---\ref{sec:lhcbimpact} and \ref{sec:atlaswz}
is considerable. Most PDFs are affected at the one-sigma level and in some
cases (such as the down and charm quarks) at up to the two-sigma level. 
This is confirmed by direct comparison of the PDFs, see
Fig.~\ref{fig:31-nnlo-lhc}. The difference between the two fits
appears to be mostly driven by the CHORUS, BCDMS and fixed-target
Drell-Yan data, whose
$\chi^2$ improves respectively by 84, 32  and 38 units when removing the
LHC data. Other datsets display much more smaller differences,
typically compatible with statistical fluctuations, and in some cases (such as
for the SLAC data) the fit quality is actually somewhat better in the
global fit. 


On the other hand, it is clear that the
shifts between PDFs without LHC data and those including them are 
compatible with the respective PDF uncertainties, and that the
uncertainties on the PDFs determined without LHC data are not so large
as to render them useless for phenomenology. We conclude that  the
default set remains considerably more accurate and should be used for
precision phenomenology. However, the use of 
PDFs determined without some or all the LHC
data may be mandatory in searches for new physics, in order to make
sure that possible new physics effects are not reabsorbed in the
PDFs. In such circumstances, we conclude that even though the uncertainty in
the PDFs without LHC data is not competitive, the level of 
deterioration is not so great as to make searches for new physics
altogether impossible. 


%%%%%%%%%%%%%%%%%%%%%%%%%%%%%%%%%%%%%%%%%%%%%%%%%%%%%%%%%%%%%%%%%%%%%
\begin{figure}[t]
\begin{center}
  \includegraphics[scale=0.38]{images/plots/xu-31-nnlo-LHC.pdf}
  \includegraphics[scale=0.38]{images/plots/xd-31-nnlo-LHC.pdf}
  \includegraphics[scale=0.38]{images/plots/xc-31-nnlo-LHC.pdf}
  \includegraphics[scale=0.38]{images/plots/xg-31-nnlo-LHC.pdf}
  \caption{\small 
Same as Fig.~\ref{fig:31-nnlo-ZpT} 
but now excluding all LHC data. Results are
shown for the up (top left), down (top right), charm (bottom left) and
gluon (bottom right) PDFs.
    %
\label{fig:31-nnlo-lhc}}
\end{center}
\end{figure}
%%%%%%%%%%%%%%%%%%%%%%%%%%%%%%%%%%%%%%%%%%%%%%%%%%%%%%%%%%%%%%%%%%%%%%

\subsection{Nuclear targets and nuclear corrections.}
\label{sec:nonucl}
The NNPDF3.1 dataset includes several measurements taken upon nuclear
targets. DIS data from the SLAC, BCDMS and NMC experiments along with
the E886 fixed-target Drell-Yan data involve measurements of deuterium.
All neutrino data and the fixed-target E605 Drell-Yan data, are obtained with heavy
nuclear targets. All of these data were already included in previous
PDF determinations, including NNPDF3.0. The impact of nuclear
corrections was studied in 
Ref.~\cite{Ball:2014uwa} and found to be under control. 
However, the much wider dataset might now permit the removal of 
these data from the global dataset: whereas removing data inevitably
entails some loss of precision, this might be more
than compensated by the increase in accuracy due to the complete elimination 
of any dependence on uncertain nuclear corrections.

In order to assess this, we performed two additional PDF
determinations 
with the NNPDF3.1
methodology. Firstly, by removing all heavy nuclear target data but keeping deuterium
data, and secondly removing all nuclear data and only keeping proton data.
The distances between the default and these two PDF sets are shown in
Fig.~\ref{fig:distances_proton}. At large $x$ the impact of
nuclear target data is significant, at the one to two sigma level,
mostly on the flavor separation of the sea. The deuterium data also have a
significant impact, particularly in the intermediate $x$ range.

 %%%%%%%%%%%%%%%%%%%%%%%%%%%%%%%%%%%%%%%%%%%%%%%%%%%%%%%%%%%%%%%%%%%%%
\begin{figure}[t]
\begin{center}
  \includegraphics[scale=1]{images/plots/distances_proton1.pdf}\\
   \includegraphics[scale=1]{images/plots/distances_proton2.pdf}
  \caption{\small Same as Fig.~\ref{fig:distances_noZpT} but now
    excluding all data with heavy nuclear targets, but keeping
    deuterium data (top) or excluding all data with any nuclear target
    and only keeping proton data (bottom) 
    \label{fig:distances_proton}
  }
\end{center}
\end{figure}
%%%%%%%%%%%%%%%%%%%%%%%%%%%%%%%%%%%%%%%%%%%%%%%%%%%%%%%%%%%%%%%%%%%%%%

%%%%%%%%%%%%%%%%%%%%%%%%%%%%%%%%%%%%%%%%%%%%%%%%%%%%%%%%%%%%%%%%%%%%%
\begin{figure}[t]
\begin{center}
  \includegraphics[scale=0.35]{images/plots/xg-31-nnlo-proton.pdf}
  \includegraphics[scale=0.35]{images/plots/xu-31-nnlo-proton.pdf}
  \includegraphics[scale=0.35]{images/plots/xd-31-nnlo-proton.pdf}
  \includegraphics[scale=0.35]{images/plots/xdbar-31-nnlo-proton.pdf}
  \caption{\small 
Same as Fig.~\ref{fig:31-nnlo-ZpT} 
but now excluding all data with heavy nuclear targets, but keeping
    deuterium data, or excluding all data with any nuclear target
    and only keeping proton data. Results are
shown for the gluon (top left), up (top right), down (bottom left) and
antidown (bottom right).
    %
\label{fig:31-nnlo-proton}}
\end{center}
\end{figure}
%%%%%%%%%%%%%%%%%%%%%%%%%%%%%%%%%%%%%%%%%%%%%%%%%%%%%%%%%%%%%%%%%%%%%%

A direct comparison of PDFs, in Fig.~\ref{fig:31-nnlo-proton}, and 
their uncertainties, in Fig.~\ref{fig:ERR-31-nnlo-datasetvar}, shows
that indeed PDFs determined with no heavy nuclear target data are
reasonably compatible with the global set, though with rather larger
uncertainties, especially for strangeness. Indeed,
best-fit results without heavy nuclear targets, or
even without deuterium data, are all compatible within their respective
uncertainties, which is consistent with the previous conclusion that
the absence of nuclear corrections for these data does not lead
to significant bias at the level of current PDF uncertainties. 
On the other hand, PDFs determined with only
proton data  while compatible to within one sigma with the
global set within their larger uncertainties, show a substantial loss of
precision. This is particularly notable for down quarks, due to the
importance of 
deuterium data in pinning down the isospin triplet PDF combinations. 

Because  deuterium data have a significant impact on the fit, one may
worry that nuclear corrections to the deuterium data are now no longer
negligible, at the accuracy of the present PDF determination. In order
to investigate this issue in greater detail, we have performed a
variant of the NNPDF3.1 NNLO default PDF determination
%, as well as of the
%NNPDF3.1 determination in which data with heavy nuclear targets have
%been removed, 
in which all deuterium data are corrected using the same
nuclear corrections as used by MMHT14 (specifically, Eqs.~(9,10) of
Ref.~\cite{Harland-Lang:2014zoa}).  


 %%%%%%%%%%%%%%%%%%%%%%%%%%%%%%%%%%%%%%%%%%%%%%%%%%%%%%%%%%%%%%%%%%%%%
\begin{figure}[t]
\begin{center}
  \includegraphics[scale=1]{images/plots/distances_global_deutcorr.pdf}
  \caption{\small Same as Fig.~\ref{fig:distances_noZpT} but now
comparing the default NNPDF3.1 to a version in which all
deuterium data have been corrected using the nuclear corrections from
Ref.~\cite{Harland-Lang:2014zoa}.
    \label{fig:distances_nuclcorr}
  }
\end{center}
\end{figure}
%%%%%%%%%%%%%%%%%%%%%%%%%%%%%%%%%%%%%%%%%%%%%%%%%%%%%%%%%%%%%%%%%%%%%%


%%%%%%%%%%%%%%%%%%%%%%%%%%%%%%%%%%%%%%%%%%%%%%%%%%%%%%%%%%%%%%%%%%%%%
\begin{figure}[t]
\begin{center}
  \includegraphics[scale=0.35]{images/plots/xu-31-nnlo-deuteron.pdf}
  \includegraphics[scale=0.35]{images/plots/xd-31-nnlo-deuteron.pdf}
  \includegraphics[scale=0.35]{images/plots/xu-ERR-31-nnlo-deuteron.pdf}
  \includegraphics[scale=0.35]{images/plots/xd-ERR-31-nnlo-deuteron.pdf}
  \caption{\small 
Same as Fig.~\ref{fig:31-nnlo-ZpT} 
but now comparing the default NNPDF3.1 to a version in which all
deuterium data have been corrected using the nuclear corrections from
Ref.~\cite{Harland-Lang:2014zoa}. Results are shown for the up (left=
and down (right) PDFs. The uncertainties are also shown (botton row).
    %
\label{fig:31-nnlo-nuclcorr}}
\end{center}
\end{figure}
%%%%%%%%%%%%%%%%%%%%%%%%%%%%%%%%%%%%%%%%%%%%%%%%%%%%%%%%%%%%%%%%%%%%%%
In terms of fit quality we find
that the inclusion of nuclear corrections leads to a
slight deterioration in the quality of the fit, with a value 
of $\chi^2/N_{\rm  dat}=1.156$, to be compared to the defaut 
$\chi^2/N_{\rm  dat}=1.148$ (see Table~\ref{tab:chi2tab_31-nlo-nnlo-30}).  
In
particular %for the global fit 
we find that for the NMC, SLAC, and
BCDMS data the values of $\chi^2/N_{\rm  dat}$ with (without) nuclear
corrections are respectively 0.94(0.95), 0.71(0.70), and 1.11(1.11).
Therefore, the addition of deuterium corrections has no significant
impact on the fit quality to these data. 

The distances between PDFs determined including deuterium corrections
and the default are shown in Fig.~\ref{fig:distances_nuclcorr}. They
are seen to be moderate and always below the half-sigma level, and confined mostly to the up and down PDFs, as expected. These PDFs  
are  shown in Fig.~\ref{fig:31-nnlo-nuclcorr}, which confirms the
moderate effect of the deuterium correction.
It should be noticed that the PDF uncertainty, also
shown in Fig.~\ref{fig:31-nnlo-nuclcorr}, is somewhat increased when
the deuterium corrections are included. The relative shift for other
PDFs are yet smaller since they are affected by larger uncertainties,
which are also somewhat increased by the inclusion of the nuclear corrections.
%In Fig.~\ref{fig:31-nnlo-nuclcorr-nonuc} the corresponding PDF
%comparison is shown for the PDF determination in which all data with
%heavy nuclear targets are removed. The conclusion is very similar,
%with the relative impact of deuterium corrections now somewhat smaller
%in comparison to the PDF uncertainties which are now larger.



%%%%%%%%%%%%%%%%%%%%%%%%%%%%%%%%%%%%%%%%%%%%%%%%%%%%%%%%%%%%%%%%%%%%%
%\begin{figure}[t]
%\begin{center}
%  \includegraphics[scale=0.35]{images/plots/xu-31-nnlo-deuteron-nonuclear.pdf}
%  \includegraphics[scale=0.35]{images/plots/xd-31-nnlo-deuteron-nonuclear.pdf}
%  \includegraphics[scale=0.35]{images/plots/xu-ERR-31-nnlo-deuteron-nonuclear.pdf}
%  \includegraphics[scale=0.35]{images/plots/xd-ERR-31-nnlo-deuteron-nonuclear.pdf}
%  \caption{\small 
%Same as Fig.~\ref{fig:31-nnlo-nuclcorr} 
%but now comparing the versions with and without deuterium corrections
%of the PDF determination in which all heavy nuclear target data have
%been removed from the   NNPDF3.1 dataset.
%    %
%\label{fig:31-nnlo-nuclcorr-nonuc}}
%\end{center}
%\end{figure}
%%%%%%%%%%%%%%%%%%%%%%%%%%%%%%%%%%%%%%%%%%%%%%%%%%%%%%%%%%%%%%%%%%%%%%%

In view of the theoretical uncertainty involved in estimating nuclear
corrections, and bearing in mind that we see no evidence of an
improvement in fit quality while we note a slight increase in PDF
uncertainties when including deuterium corrections using the model of
Ref.~\cite{Harland-Lang:2014zoa}, we conclude that the impact of
deuterium corrections on the NNPDF3.1 results is sufficiently small that they
may be safely ignored even within the current high precision of PDF
determination. Nevertheless, more detailed dedicated studies of
nuclear corrections, also in relation to the construction of nuclear
PDF sets, may well be worth pursuing in future studies.

In conclusion, for the time being it is still appears advantageous to retain
nuclear target data in the global dataset for general-purpose PDF
determination. However, if very high accuracy is required (such as,
for instance, in the determination of standard model parameters) it
might be preferable to use PDF sets from which all data with nuclear
targets have been omitted. 

%%%%%%%%%%%%%%%%%%%%%%%%%%%%%%%%%%%%%%%%%%%%%%%%%%%%%%%%%%%%%%%%%%%%%
\begin{figure}[t]
  \begin{center}
    \includegraphics[scale=0.32]{images/plots/xu-ERR-31-nnlo-datasetvar.pdf}
    \includegraphics[scale=0.32]{images/plots/xubar-ERR-31-nnlo-datasetvar.pdf}
    \includegraphics[scale=0.32]{images/plots/xd-ERR-31-nnlo-datasetvar.pdf}
    \includegraphics[scale=0.32]{images/plots/xdbar-ERR-31-nnlo-datasetvar.pdf}
  \includegraphics[scale=0.32]{images/plots/xs-ERR-31-nnlo-datasetvar.pdf}
  \includegraphics[scale=0.32]{images/plots/xsbar-ERR-31-nnlo-datasetvar.pdf}
  \includegraphics[scale=0.32]{images/plots/xc-ERR-31-nnlo-datasetvar.pdf}
  \includegraphics[scale=0.32]{images/plots/xg-ERR-31-nnlo-datasetvar.pdf}
  \caption{\small Comparison of the relative
    PDF uncertainties at $Q=100$ between the NNPDF3.1 and the no heavy nuclei, 
proton--only and collider--only PDF determinations. The uncertainties shown are all
    normalized to the NNPDF3.1 central value.
    \label{fig:ERR-31-nnlo-datasetvar}
  }
\end{center}
\end{figure}
%%%%%%%%%%%%%%%%%%%%%%%%%%%%%%%%%%%%%%%%%%%%%%%%%%%%%%%%%%%%%%%%%%%%%%

\subsection{Collider-only parton distributions}
\label{sec:collideronly}

A yet more conservative option to that discussed in the previous
Section is to retain only collider data from HERA, the
Tevatron and the LHC. The motivation for this suggestion,
first presented in the NNPDF2.3 study~\cite{Ball:2012cx}, is that
this excludes data taken at low
scales, which may be subject to potentially large perturbative and non-perturbative
corrections. Furthermore, data taken on nuclear targets, and all of the older datasets 
are eliminated, thereby leading to a more reliable set of PDFs.
However, previous collider-only PDFs had very large uncertainties, due to the 
limited collider dataset then available.

In order to re-assess the situation with the current, much wider LHC dataset,
we have repeated a collider-only PDF determination. This amounts to
repeating the proton only PDF determination 
described in the previous Section, but now with the proton fixed target data also removed. 
The distances between the ensuing
PDFs and the default NNPDF3.1 are shown in
Fig.~\ref{fig:distances_collider}. Comparing to
Fig.~\ref{fig:distances_proton} we notice
that distances are similar for most PDFs, the main exception
being the gluon, which in the intermediate $x$ region has now shifted 
by almost two sigma. 

%%%%%%%%%%%%%%%%%%%%%%%%%%%%%%%%%%%%%%%%%%%%%%%%%%%%%%%%%%%%%%%%%%%%%
\begin{figure}[t]
\begin{center}
  \includegraphics[scale=1]{images/plots/distances_nnpdf31_collider.pdf}
  \caption{\small 
Same as Fig.~\ref{fig:distances_noZpT} but now
 only keeping collider data.
    \label{fig:distances_collider}
  }
\end{center}
\end{figure}
%%%%%%%%%%%%%%%%%%%%%%%%%%%%%%%%%%%%%%%%%%%%%%%%%%%%%%%%%%%%%%%%%%%%%%

This is confirmed by a direct comparison of PDFs in
Fig.~\ref{fig:pdfs-collideronly} and their uncertainties in Fig.~\ref{fig:ERR-31-nnlo-datasetvar}. 
The valence quarks (especially up) are reasonably stable, but the sea now is quite unstable upon 
removal of all the fixed target data, visibly more than in the proton only set
Fig.~\ref{fig:31-nnlo-proton}, with in particular a substantial increase in the 
uncertainty of the anti-up quark distribution at large $x$. Furthermore, the gluon, which in
Fig.~\ref{fig:31-nnlo-proton} was quite stable in the proton-only PDF
set, now undergoes a significant downward shift at intermediate $x$, even though its 
uncertainty is not substantially increased. 

We conclude that, despite impressive improvements due to recent LHC measurements, 
a collider-only PDF determination is still not very useful for general phenomenological applications.


%%%%%%%%%%%%%%%%%%%%%%%%%%%%%%%%%%%%%%%%%%%%%%%%%%%%%%%%%%%%%%%%%%%%%
\begin{figure}[t]
\begin{center}
  \includegraphics[scale=0.38]{images/plots/xg-31-nnlo-collider.pdf}
  \includegraphics[scale=0.38]{images/plots/xu-31-nnlo-collider.pdf}
  \includegraphics[scale=0.38]{images/plots/xd-31-nnlo-collider.pdf}
  \includegraphics[scale=0.38]{images/plots/xdbar-31-nnlo-collider.pdf}
  \caption{\small 
Same as Fig.~\ref{fig:31-nnlo-ZpT} 
but now only keeping collider data.
Results are
shown for the gluon (top left), up (top right), down (bottom left) and
antidown (bottom right).
    \label{fig:pdfs-collideronly}
  }
\end{center}
\end{figure}
%%%%%%%%%%%%%%%%%%%%%%%%%%%%%%%%%%%%%%%%%%%%%%%%%%%%%%%%%%%%%%%%%%%%%





